\section{Kopplingen}
\label{sec:wiring}

\subsection{SPI mellan korten}

Fyra ledningar plus matning och jord, enligt kontraktets pinnkonfiguration. Två saker är värda att
säga ut:

\textbf{\code{VDDIO2} först.} \code{PORTC} är MVIO-domänen och matas därifrån. Utan den kopplingen
gör porten ingenting alls, och symptomet ser ut precis som en trasig SPI-implementation.

\textbf{Gemensam jord är inte valfritt.} Två kort utan gemensam jord kommunicerar inte, och
symptomet är obegripligt.

\subsection{CAN-sidan}

FPGA-nodens CAN-sida är en \textbf{enledarbuss med öppen dränering} på tre GPIO-pinnar:
\code{tx_bus} (vad noden driver), \code{bus_en} (om den driver alls) och \code{rx_bus} (bussens
nivå in).

För att en bussanalysator ska kunna se trafiken måste den enledarbussen bli en riktig
\textbf{differentiell} CAN-buss. Det är transceiverns jobb, och dess sändaringång ska vara
\textbf{låg för en dominant bit}:

\begin{narrowtable}{@{}llll@{}}
\thead{\code{bus_en}} & \thead{\code{tx_bus}} & \thead{Noden} & \thead{Transceiverns ingång} \\ \midrule
0 & --- & driver inte & 1 (recessiv) \\
1 & 0 & dominant & 0 \\
1 & 1 & recessiv & 1 \\
\end{narrowtable}

Hur de två signalerna slås ihop till den enda ledningen är hårdvarusidans del av kopplingen; stäm
av det innan du kopplar. En nod vars \code{bus_en} ignorerats driver bussen hela tiden och
blockerar alla andra.

\textbf{Terminering: 120~\textohm{} i vardera änden}, inte i mitten och inte fler än två. En
felterminerad buss fungerar ofta ändå på en kort kabel, och slutar fungera precis när du börjat
lita på den.

\section{Stegen}
\label{sec:ladder}

Klättra dem i ordning. Varje steg har ett eget kvitto, och poängen är att \textbf{aldrig ha mer än
en okänd sak åt gången}. Hoppa över ett steg och du felsöker fyra lager samtidigt.

\begin{booktable}{@{}llL@{}}
\thead{Steg} & \thead{Gör} & \thead{Kvitto} \\ \midrule
0 & Matning och jord & MVIO-statusen rapporterar att \code{VDDIO2} är inom intervallet. \\
1 & \code{MOSI} till \code{MISO}, utan FPGA & Du får tillbaka exakt den byte du skickade. Prova
  \code{0x00}, \code{0xFF}, \code{0xAA}, \code{0x55}. \\
2 & Koppla in FPGA; skriv och läs \code{TX_ID} & \code{0xFFFFFFFF} läses tillbaka som
  \code{0x000007FF}. \\
3 & Läs \code{STATUS} upprepat & \code{0x00000001} --- TX klar satt, övriga låga. \\
4 & Koppla in transceivern; kör en \code{send()} & Aktivitet på bussen; \code{STATUS} bit 0 går
  låg och kommer tillbaka. \\
5 & Låt en andra nod sända & \code{receive()} returnerar \code{true} med rätt innehåll. \\
\end{booktable}

Två av stegen förtjänar en kommentar.

\textbf{Steg 1 är \code{AvrSpiTransport}s enda test.} Det är den enda klassen i projektet som
inget automattest täcker, eftersom den pratar med ett fysiskt register. Fungerar loopbacken inte
ligger felet i din kod eller i konfigurationen --- FPGA:n är inte inkopplad och kan inte vara
orsaken.

\textbf{Steg 2 är det avgörande.} Maskeringen till elva bitar bevisar att du pratar med
\emph{registerbanken} och inte med ett eko av din egen ledning. En loopback ger tillbaka det du
skickade; en registerbank ger tillbaka det den lagrade.

\section{Felsökning}
\label{sec:troubleshooting}

\begin{booktable}{@{}lLl@{}}
\thead{Symptom} & \thead{Trolig orsak} & \thead{Kontrollera} \\ \midrule
Inget alls, inga klockpulser & Pinnmuxen inte satt & Steg 1 \\
Porten gör ingenting & \code{VDDIO2} inte matad & MVIO-statusen \\
Alla läsningar ger \code{0x00000000} & \code{MISO} inte kopplad, eller fel pinne & Byt
  \code{MOSI}/\code{MISO} och se om det ändras \\
Alla läsningar ger \code{0xFFFFFFFF} & \code{MISO} flytande & Kopplingen \\
Läsningarna ligger en byte fel & Dataregistret läses inte i \code{transfer()} & \kapref{ch:avr} \\
Rätt för små tal, fel för stora & Teckenutvidgning i \code{readReg()} & \kapref{ch:byte} \\
Andra transaktionen ignoreras & Två transaktioner under en låg \code{SS} & \kapref{ch:spi} \\
Framen kommer fram bakvänd & \code{TX_DATA_LO}/\code{HI} förväxlade & \kapref{ch:driver} \\
Samma frame om och om igen & \code{RX_ACK} skrivs inte & \kapref{ch:driver} \\
\code{hasError()} alltid \code{true} & \code{clearError()} skriver \code{0x1} & \kapref{ch:driver} \\
Bussen konstant dominant & En nods \code{bus_en} ignorerad & Kopplingstabellen \\
Fungerar på kort kabel, inte lång & Terminering & 120~\textohm{} i vardera änden \\
\end{booktable}

\subsection{Arbetsgången}

\begin{enumerate}
  \item Vilket steg var det sista som gav rätt kvitto?
  \item Vad är den \textbf{enda} saken som tillkommit sedan dess?
  \item Formulera en hypotes om vad som är fel.
  \item Bestäm vilken mätning som skulle avgöra saken.
  \item Gör mätningen.
\end{enumerate}

\begin{aside}
\note Att byta sladdar tills det fungerar är inte felsökning, och det tar längre tid. Skillnaden
mellan de två arbetssätten är ungefär en eftermiddag per fel, och den är den mest överförbara
färdigheten i hela kursen.
\end{aside}

\subsection{Att felsöka någon annans halva}

En bugg under bring-upen kan ligga i din kod, i den andra gruppens VHDL, eller i en kopplingstråd.
Tre regler gör det hanterbart:

\begin{itemize}
  \item \textbf{Misstänk din egen halva först.} Inte av ödmjukhet, utan för att du kan ändra den
    och mäta på den.
  \item \textbf{Bevisa innan du rapporterar.} "Det fungerar inte" är inte en felrapport.
    "Steg 2 ger \code{0x000007FF} för \code{TX_ID} men \code{0x00000000} för \code{TX_DLC}" är det.
  \item \textbf{Ändra en sak i taget, och skriv ned vad.} Två personer som felsöker samma
    koppling samtidigt gör varandras mätningar meningslösa.
\end{itemize}
